% =====================================================================
%  CHAPTER 3: 青春期发育（对应大纲第三章）
%  内容来源：往届笔记第五章，参考PPT第五章
% =====================================================================
\chapter{青春期发育}
\setcounter{chapter}{3}
\chapterintro{
\textbf{掌握：}青春期发动机理，青春期生长发育的特点和卫生问题。\newline
\textbf{熟悉：}青春期基本概念，青春期内分泌。\newline
\textbf{其他：}青春期性发育的特殊表现。
}

% ===== 第一部分：掌握内容 =====
\section{青春期发动机理}\difficulty

\subsection{神经-内分泌调控机制}

\begin{keybox}
\textbf{青春期发育的启动机制：}

青春期内分泌调节的核心是\textbf{神经-内分泌调控机制}。青春期发育的启动机制包括：

\begin{enumerate}[leftmargin=*]
    \item \textbf{中枢神经抑制机制减弱}——儿童期中枢神经系统对下丘脑的抑制作用逐渐减弱。
    \item \textbf{HPG轴对性激素负反馈敏感性下降}——下丘脑-垂体-性腺（HPG）轴对性激素负反馈的敏感性降低，导致GnRH脉冲式分泌增加。
\end{enumerate}
\end{keybox}

\subsection{青春期内分泌变化}

\begin{keybox}
\textbf{1. 性腺功能初现：}由下丘脑分泌的促性腺激素释放激素（GnRH）发动。男童：睾丸体积增大；女童：乳房发育。

\textbf{2. 肾上腺皮质功能初现：}由肾上腺分泌的肾上腺雄激素发动。出现时间较早（6-8岁），肾上腺雄激素分泌开始程序性增加，此时雄激素活性比睾酮作用弱，对男童作用更弱、显现不出作用；但对女童第二性征发育作用明显。
\end{keybox}

\section{青春期生长发育的特点}

\subsection{青春期体格发育特点}

\begin{keybox}
\textbf{1. 青春期生长突增}

（1）\textbf{生长突增（growth spurt）：}进入青春期后最引人注目的形态变化是儿童少年体格生长出现的突发性快速生长现象。

①\textbf{身高速度高峰（peak height velocity, PHV）：}以身高为代表，生长速度在童年期较平稳的每年4-5cm基础上出现快速增长，1-2年后达到高峰。PHV发生时的年龄称\textbf{身高速度高峰年龄（age at peak height, PHA）}。

②\textbf{体重速度高峰（peak weight velocity, PWV）}和\textbf{体重速度高峰年龄（age at peak weight, PWA）}：儿童少年体格发育的另一重要指标，在青春期发育阶段的生长发育特点与身高发育特点基本一致。

\textbf{2. 青春期生长突增的速度变化}

（1）年龄别生长速度曲线和按PHA生长速度曲线存在明显不同：
\begin{itemize}
    \item 由于青春期体格生长突增高峰年龄的个体变异突出，按年龄别生长速度曲线中儿童身高生长突增高峰值明显小于按PHA生长速度曲线中儿童身高生长突增高峰值。
    \item 因此，评价个体儿童青春期生长发育速度时，不能单独应用群体按年龄生长速度参考值，可将按年龄生长速度参考值和按PHA生长速度参考值结合考虑。
\end{itemize}

（2）男女身高、体重生长曲线出现\textbf{两次交叉现象}：
\begin{itemize}
    \item \textbf{第一次交叉：}女性因体格生长突增开始早，青春期开始早期女性体格生长水平超过同龄男性。
    \item \textbf{第二次交叉：}随年龄增长女性体格生长逐步进入缓慢阶段，而男性体格生长突增开始，在男性体格生长速度达到高峰前，其生长水平已超过女性。
\end{itemize}

（3）青春期身体各部分的生长速度：青春前期，上下肢增长早于躯干，下肢增长稍早于上肢。顺序大致为：足长→下肢长→手长→上肢长。足长生长首先加速，也最早停止（一般14-15岁后）；可根据青少年足长最先突增又较早停止生长的特点，用足长预测成年身高。

\textbf{3. 青春期生长发育类型}

（1）\textbf{3种成熟类型生长特点}

\begin{table}[H]
\centering
\caption{三种成熟类型的生长特点比较}
\begin{tabularx}{\textwidth}{l>{\raggedright\arraybackslash}p{4cm}>{\raggedright\arraybackslash}p{4cm}>{\raggedright\arraybackslash}p{4cm}}
\hline
\textbf{特征} & \textbf{早熟型} & \textbf{一般型} & \textbf{晚熟型} \\
\hline
体重/身高比值 & 一直较高 & 中间 & 一直较低 \\
\hline
体型特征 & 盆宽、窄肩，矮胖 & 中间 & 盆窄、肩宽，瘦高 \\
\hline
青春期启动 & 女孩8-9岁，男孩10-11岁 & 中间 & 女孩10-11岁，男孩12-13岁 \\
\hline
生长期 & 较短，约1年 & 约2年 & 较长，可达2-3年甚至3年 \\
\hline
性别偏向 & 女性体型特征 & —— & 男性体型特征 \\
\hline
\end{tabularx}
\end{table}

（2）\textbf{成熟类型与生长突增}
\begin{itemize}
    \item 模式I：生长突增起始于平均年龄，成年身高位于平均水平。
    \item 模式II：生长突增发生较早，突增结束年龄也较早，生长期较短，成年身高低于平均水平。
    \item 模式III：儿童期和青春早期生长低于同龄者，但较晚的生长突增和较长的生长期导致其成年身高达到甚至超过平均水平。
    \item 模式IV：遗传性（家族性）高身材。
    \item 模式V：遗传性（家族性）矮身材。
\end{itemize}
\end{keybox}

\subsection{青春期性发育特点}

\begin{keybox}
\textbf{1. 男性性发育}

（1）\textbf{性器官形态发育：}各指征的出现顺序大致相似：睾丸→阴茎、身高突增（一年后）。

（2）\textbf{第二性征发育：}阴毛→腋毛→胡须→变声、喉结出现。绝大多数男孩在18岁前完成第二性征发育。

（3）\textbf{性功能发育：}遗精是男性生殖功能开始发育成熟的重要标志之一。首次遗精一般发生于12-18岁间。

\textbf{2. 女性性发育}

（1）\textbf{性器官形态发育：}进入青春期后，在垂体/性腺分泌的促卵泡激素、黄体生成素、促性腺激素作用下，内外生殖器迅速发育。

（2）\textbf{第二性征发育：}乳房→阴毛→腋毛发育。乳房发育平均开始于11岁，历时约4年。身高生长突增几乎与乳房发育同时或稍早开始，而身高突增高峰出现年龄通常在乳房发育后1年左右。

（3）\textbf{性功能发育：}月经初潮（在身高突增高峰后1年左右）。
\end{keybox}

\begin{exambox}
\textbf{△ 关键标志：}
\begin{itemize}
    \item 进入青春期第一信号：身高突增
    \item 性发育第一信号：男孩——睾丸增大，女孩——乳房发育
    \item 生殖功能开始成熟的标志：男孩——首次遗精，女孩——月经初潮
    \item 第二性征：人或动物发育后表现出的与生殖系统无直接关系、而与性别有关的机体外表特征
\end{itemize}
\end{exambox}

\subsection{青春期卫生问题}

\begin{tipbox}
\textbf{青春期主要卫生问题：}

\begin{enumerate}[leftmargin=*]
    \item \textbf{青春发动时相异常}——提前（early）、推迟（delay）均对身心健康有不利影响。一般以人群青春发动期生殖系统发育各种事件初现时间的前第25百分位数（P25）作为青春发动期时相提前的划界值。
    \item \textbf{青春期心理卫生问题}——青春期是性心理障碍开始表现的关键期，出现异性疏远期→异性爱慕期→两性初恋期的发展阶段，以及性生理发育成熟与性心理相对幼稚的矛盾等。
    \item \textbf{营养卫生问题}——生长突增期对营养需求增加，易出现营养不足或过剩。
    \item \textbf{伤害与暴力}——青春期冒险行为增加，伤害成为主要健康威胁之一。
    \item \textbf{健康危险行为}——吸烟、饮酒、不当饮食行为、缺乏体力活动等。
\end{enumerate}
\end{tipbox}

% ===== 第二部分：熟悉内容 =====
\section{青春期基本概念}

\begin{defbox}
\textbf{青春期（adolescence）（WHO）：}个体从出现第二性征到性成熟的生理发育过程，是个体从儿童认知方式发展到成人认知方式的心理过程，是个体从经济的依赖性到相对独立状态的过渡，10-19岁。

\textbf{青春发育进程（pubertal procession）：}描述某一时点，某个体在整个青春期发育过程中所处的绝对位置。特点：
\begin{enumerate}[leftmargin=*]
    \item 体格生长加速，以身高为代表的体格指标出现第二次生长突增。
    \item 各内脏器官体积增大、重量增加，功能日渐成熟。
    \item 内分泌功能活跃，与生长发育有关的激素分泌明显增加。
    \item 生殖系统功能发育骤然增快并迅速成熟，到青春晚期已具备繁殖后代的能力。
    \item 男女外生殖器和第二性征发育，使男女两性外部形态特征差别更明显。
    \item 在体格、功能发育同时，青春期心理发展加速，产生相应心理、行为变化，出现很多青春期特有的心理-行为问题。
\end{enumerate}

\textbf{青春期分期：}早、中、晚三个阶段。
\begin{enumerate}[leftmargin=*]
    \item \textbf{青春早期：}主要表现为生长突增，出现身高突增高峰，性发育开始，一般持续约2-3年。
    \item \textbf{青春中期：}以性器官、第二性征迅速发育为特征，出现月经初潮（女）或首次遗精（男），持续2-3年。
    \item \textbf{青春后期：}体格生长速度明显减慢，但仍有增长，直至骨骺完全融合，性器官和第二性征持续发育至成人水平，社会心理发展加速，通常持续2年左右。
\end{enumerate}

\textbf{青春发动期（puberty）：}青春期生殖系统发育与成熟的生物学变化过程，是青春期的早期阶段，时间跨度平均约4.5年（1.5-6年）。个体生理变化明显，包括体格生长突增、性腺和生殖器发育、第二性征出现等。

\textbf{青春发动时相（puberty timing）：}在青春发动期，生殖系统发育的各种事件按特定时间模式进行，从生长突增、乳房发育、腋毛生长、阴毛生长，到月经初潮等，这种青春发动期生殖系统发育的各种事件初现的时间模式。
\end{defbox}

\section{青春期内分泌}

\begin{keybox}
\textbf{影响青春期发育的主要内分泌激素：}

\textbf{1. 下丘脑}
\begin{itemize}
    \item 释放激素：促生长素释放激素、促甲状腺激素释放激素、促肾上腺皮质激素释放激素、促卵泡激素释放激素、黄体生成素释放激素、催乳素释放因子和促黑激素释放激素。
    \item 释放抑制激素：生长抑素、催乳素释放抑制因子和促黑激素释放抑制因子。
\end{itemize}

\textbf{2. 腺垂体}
\begin{itemize}
    \item \textbf{生长素（GH）：}促使氨基酸进入细胞加速蛋白质合成；分解脂肪使游离脂肪酸增加；抑制葡萄糖氧化减少糖原消耗。受下丘脑GHRH与GHRIH双重调节。生长素分泌与睡眠有关，入睡时分泌明显增加，入睡后60min左右血中浓度达到高峰；慢波睡眠时相分泌增加，快波睡眠时相分泌减少。
    \item \textbf{催乳素（PRL）：}促进乳腺发育，启动和维持乳腺泌乳。受下丘脑双重调节，正常生理情况下抑制占优势，所以青春期女性乳腺发育成熟但无乳汁生成。
    \item \textbf{促激素：}促甲状腺激素（TSH）、促肾上腺皮质激素（ATCH）、促卵泡激素（FSH）（促进卵泡成熟）和黄体生成素（LH）（促进排卵）。
\end{itemize}

\textbf{3. 性腺}
\begin{itemize}
    \item \textbf{睾丸-雄激素：}睾酮（最重要）、双氢睾酮和脱氢表雄酮。睾酮功能：胚胎发育期对男性胎儿生殖器的分化起关键作用；儿童期促进体内蛋白质合成和骨骼、肌肉发育。
    \item \textbf{卵巢-雌激素：}促进女性青春期乳腺发育最关键的激素。
\end{itemize}

\textbf{4. 甲状腺：}甲状腺素（TH），功能：调节机体物质代谢、促进生长发育、维持神经和心血管功能。出生后功能低下→呆小症/克汀病；缺碘→地方性克汀病。

\textbf{5. 肾上腺：}肾上腺皮质与青春期生长发育有关。雄激素（脱氢表雄酮）：调节性器官和第二性征发育，对女性阴毛、腋毛等促进较明显，对男性影响较小。

\textbf{6. 胰岛：}胰岛素与生长激素有协同作用，可与IGF-I受体结合而刺激生长。

\textbf{7. 松果体：}褪黑素，功能：抑制哺乳动物生殖系统发育。
\end{keybox}

% ===== 第三部分：其他内容 =====
\section{青春期性发育的特殊表现}

\subsection{青春期性心理发育}

\begin{tipbox}
\textbf{青春期性心理发展阶段：}

\begin{enumerate}[leftmargin=*]
    \item \textbf{异性疏远期：}性发育早期体态剧烈变化使青少年心理出现适应困难，内心慌乱不安，表现出对性的抵触现象，甚至反感。大部分男女生开始相互疏远异性。
    \item \textbf{异性爱慕期：}随着青春期性心理逐步发展，男女都开始渴望了解异性，主动接近异性，同时开始萌发对异性的向往和憧憬。
    \item \textbf{两性初恋期：}性意识的发展逐渐趋于成熟，对异性的情感充满浪漫和幻想，对异性爱慕逐渐趋于专一。男女生情感都不稳定，并有明显占有心理。
\end{enumerate}

\textbf{青春期性心理发展矛盾现象：}
\begin{enumerate}[leftmargin=*]
    \item 性生理发育成熟与性心理相对幼稚的矛盾。
    \item 自我意识迅猛发展与社会成熟度相对迟缓的矛盾。
    \item 情感激荡释放与外部表露趋向内隐的矛盾。
\end{enumerate}

△青春期是性心理障碍开始表现的关键期。
\end{tipbox}

% ----- 复习题 -----
\reviewcontent{
\section{复习题}

\begin{enumerate}[leftmargin=*, itemsep=8pt]
    \item \textbf{【名词解释】}青春期；青春发动期；青春发动时相；青春发育进程；生长突增；身高速度高峰（PHV）；第二性征；性腺功能初现；肾上腺皮质功能初现
    \item \textbf{【简答题】}简述青春期发育的启动机制。
    \item \textbf{【简答题】}简述青春期体格发育的主要特点。
    \item \textbf{【简答题】}比较早熟型、一般型和晚熟型三种成熟类型的生长特点。
    \item \textbf{【简答题】}简述男女身高、体重生长曲线两次交叉现象及其原因。
    \item \textbf{【简答题】}简述影响青春期发育的主要内分泌激素及其功能。
    \item \textbf{【简答题】}简述青春期性心理发展的三个阶段和矛盾现象。
\end{enumerate}

\subsection*{参考答案要点}

\begin{enumerate}[leftmargin=*, itemsep=5pt]
    \item \textbf{青春期}：个体从出现第二性征到性成熟的生理发育过程，是个体从儿童认知方式发展到成人认知方式的心理过程，是个体从经济的依赖性到相对独立状态的过渡，10-19岁。\textbf{青春发动期}：青春期生殖系统发育与成熟的生物学变化过程，是青春期的早期阶段，时间跨度平均约4.5年（1.5-6年）。\textbf{青春发动时相}：在青春发动期，生殖系统发育的各种事件按特定时间模式进行，从生长突增、乳房发育、腋毛生长、阴毛生长，到月经初潮等，这种青春发动期生殖系统发育的各种事件初现的时间模式。\textbf{青春发育进程}：描述某一时点，某个体在整个青春期发育过程中所处的绝对位置。\textbf{生长突增}：进入青春期后最引人注目的形态变化是儿童少年体格生长出现的突发性快速生长现象。\textbf{PHV}：以身高为代表，生长速度在童年期较平稳的每年4-5cm基础上出现快速增长，1-2年后达到的高峰。\textbf{第二性征}：人或动物发育后表现出的与生殖系统无直接关系、而与性别有关的机体外表特征。\textbf{性腺功能初现}：由下丘脑分泌的促性腺激素释放激素（GnRH）发动，男童睾丸体积增大，女童乳房发育。\textbf{肾上腺皮质功能初现}：由肾上腺分泌的肾上腺雄激素发动，出现时间较早（6-8岁），肾上腺雄激素分泌开始程序性增加。

    \item (1) 中枢神经抑制机制减弱；(2) HPG轴对性激素负反馈敏感性下降，导致GnRH脉冲式分泌增加。神经-内分泌调控机制是青春期内分泌调节的核心。

    \item (1) 青春期生长突增：出现PHV和PWV；(2) 男女身高体重曲线两次交叉；(3) 身体各部分生长速度不同，顺序为足长→下肢长→手长→上肢长；(4) 青春期身体各部分比例不断变化，青春早期长臂长腿体态；(5) 三种成熟类型（早熟型、一般型、晚熟型）生长特点不同。

    \item 早熟型：体重/身高比值高，盆宽窄肩矮胖，启动早（女8-9男10-11），生长期短。一般型：各特征居中，突增时间约2年。晚熟型：体重/身高比值低，盆窄肩宽瘦高，启动晚（女10-11男12-13），生长期长（2-3年+）。

    \item 第一次交叉：女性体格生长突增开始早，青春期早期女性体格水平超过同龄男性。第二次交叉：女性生长进入缓慢阶段，男性生长突增开始并超过女性。

    \item (1) 下丘脑释放/抑制激素；(2) 腺垂体：生长素（促进蛋白质合成）、催乳素（促进乳腺发育）、促激素（调节内分泌器官）；(3) 性腺：睾酮（男性生殖器分化、蛋白质合成）、雌激素（乳腺发育最关键激素）；(4) 甲状腺素（调节代谢、促进发育）；(5) 肾上腺雄激素（调节第二性征）；(6) 胰岛素（与GH协同刺激生长）；(7) 褪黑素（抑制生殖系统发育）。

    \item 三个阶段：(1)异性疏远期→(2)异性爱慕期→(3)两性初恋期。三大矛盾：性生理成熟与性心理幼稚的矛盾、自我意识发展与社会成熟度迟缓的矛盾、情感激荡释放与外部内隐的矛盾。
\end{enumerate}
}

\newpage
